%! TeX program = lualatex

\documentclass[12pt]{article}

\usepackage{cmap}

\usepackage[english, russian]{babel}

\usepackage{float}
\usepackage{fontspec}
\usepackage{libertinus}
\setmonofont{Source Code Pro}

\usepackage{microtype}
\usepackage{graphicx}
\usepackage{xcolor} % Цвета

\usepackage{subfig}

\usepackage{listings}

\usepackage[left=2.5cm, right=2.5cm, top=2cm, bottom=2cm]{geometry}

\usepackage[mocha, textcolor=true, pagecolor=true]{catppuccinpalette}
% \usepackage[latte, textcolor=true, pagecolor=true]{catppuccinpalette}


\definecolor{codegreen}{rgb}{0,0.6,0}
\definecolor{codegray}{rgb}{0.5,0.5,0.5}
\definecolor{codepurple}{rgb}{0.796, 0.651, 0.969}
\definecolor{codeBlue}{rgb}{0.537, 0.706, 0.980}

\lstdefinestyle{mystyle}{
    commentstyle=\color{codegreen},
    keywordstyle=\color{codeBlue},
    numberstyle=\tiny\color{codegray},
    stringstyle=\color{codepurple},
    basicstyle=\ttfamily\footnotesize,
    breakatwhitespace=false,         
    frame=leftline,
    xleftmargin=8pt,
    breaklines=true,                 
    captionpos=b,                    
    keepspaces=true,                 
    numbers=left,                    
    numbersep=5pt,                  
    showspaces=false,                
    showstringspaces=false,
    showtabs=false,                  
    tabsize=2
}

\lstset{style=mystyle}
\renewcommand{\lstlistingname}{}

\begin{document}

\begin{center}
	\section*{Лабораторная работа №1}
	Вариант №3
\end{center}

\noindent Листинг библиотеки классов (dll):

\begin{lstlisting}[language=Python]
using System;

namespace dll
{
    public class Machine
    {
        public string Metal { get; set; } = "None";
        public int Age { get; set; } = 0;
        public bool Break { get; set; } = false;

        public void Repair()
        {
            if (Break)
            {
                Break = false;
                Console.WriteLine("Car is fixed");
            }

            else
                Console.WriteLine("Car is in good condition");

        }

        public void setAge(int new_age)
        {
            Age = new_age;
            Console.WriteLine($"Set age: {Age}");
        }

        public void ShowInfo()
        {
            string Status = Checking();
            Console.WriteLine($"\nMachine age: {Age}\nStatus: {Status}\nMetal: {Metal}");
        }

        public string Checking()
        {
            if (!Break)
                return "serviceable";
            else
                return "break";
        }
    }

    public class Car : Machine
    {
        public int id { get; set; } = 0;
        public string Model { get; set; } = "None";
        public int Places = 0;
        public Engine engine = new Engine();
        public bool Drive = false;


        public void DriveToggle()
        {
            if (Drive) { Drive = false; }
            else { Drive = true; }

            if (Drive && !Break)
            {


                if (!engine.Work)
                {
                    engine.Start();
                }

                Console.WriteLine("Drive");
            }

            else if (Break)
            {
                Drive = false;
                Console.WriteLine("Car is broken!");
            }
            else
            {
                if (engine.Work)
                {
                    engine.Stop();
                }
                Console.WriteLine("Stop drive");
            }
        }

        public void Honk() { Console.WriteLine("HONK!"); }

        public void ReplaceEngine()
        {
            if (engine.Work) { engine.Stop(); }

            Console.WriteLine("Model engine: ");
            string model = Console.ReadLine();

            Console.WriteLine("Cylinders engine: ");
            int cylinders = int.Parse(Console.ReadLine());

            Console.WriteLine("HorsePower engine: ");
            int horse_power = int.Parse(Console.ReadLine());

            Engine new_engine = new Engine { model_engine = model, Cylinders = cylinders, HorsePower = horse_power };
            engine = new_engine;

        }

        public void modelEngine()
        {
            Console.WriteLine(engine.model_engine);

        }

        public Car CreateCar()
        {
            Car newCar = new Car();

            Console.WriteLine("Model: ");
            newCar.Model = Console.ReadLine();

            Console.WriteLine("Places: ");
            newCar.Places = int.Parse(Console.ReadLine());

            Console.WriteLine("--Machine Properties--");
            Console.WriteLine("Metal: ");
            newCar.Metal = Console.ReadLine();

            Console.WriteLine("Age: ");
            newCar.Age = int.Parse(Console.ReadLine());

            Console.WriteLine("Break (true/false): ");
            newCar.Break = bool.Parse(Console.ReadLine());

            Console.WriteLine("--Engine Properties--");
            /*newCar.engine = new Engine();*/
            newCar.engine.createEngine();

            return newCar;

        }

    }

    public class Engine
    {
        public int id { get; set; } = 0;
        public string model_engine { get; set; } = "None";
        public int Cylinders { get; set; } = 0;
        public int HorsePower { get; set; } = 0;
        public bool Work = false;

        public Engine createEngine()
        {

            Console.WriteLine("Model: ");
            string Model_engine = Console.ReadLine();
            Console.WriteLine("Cylinders: ");
            int cylinders = int.Parse(Console.ReadLine());
            Console.WriteLine("HorsePower: ");
            int horsePower = int.Parse(Console.ReadLine());

            return new Engine { Cylinders = Cylinders, HorsePower = horsePower, model_engine = Model_engine };
        }

        public void Start()
        {
            if (!Work) { Work = true; } else { Console.WriteLine("The engine is already running"); }
        }

        public void Stop()
        {
            if (Work) { Work = false; } else { Console.WriteLine("The engine is no longer running"); }
        }


        public void Upgrade(int extraHorsPower)
        {
            HorsePower = extraHorsPower;
        }

    }
}
\end{lstlisting}
\vspace{12pt}

\noindent Листинг консольного приложения

\begin{lstlisting}[language=Python]
using dll;

class Program
{
    static public void Main()
    {
        bool exit = false;

        do
        {
            Console.Clear();
            Console.WriteLine("1. Create car");
            Console.WriteLine("2. Info car");
            Console.WriteLine("0. Exit");
            int choice = int.Parse(Console.ReadLine());

            switch (choice)
            {
                case 1:
                    Car newCar = new Car();
                    newCar.CreateCar();
                    break;
                case 0:
                    exit = true;
                    Console.Clear();
                    break;
            }

        } while (!exit);
    }
}

\end{lstlisting}

\subsection*{Результаты работы программы:}

\end{document}
